\subsection*{Phase I-C: Higher Confinement and Septinion Unfolding}
\addcontentsline{toc}{subsection}{Phase I-C: Higher Confinement and Septinion Unfolding}

This third unfolding block turns closure stability into a higher confinement
reading and makes explicit why the septinion/seption interpretation is not an
external axiom but a structurally selected reading of the admissible closure
class.

\begin{lemma}[Higher confinement lemma]
\label{lem:higher-confinement-lemma}
If a structural class is closure-stable and remains invariant under repeated
compensator-compatible normalization, then it admits a higher confinement
reading.
\end{lemma}

\begin{proof}
Closure stability excludes uncontrolled drift out of the admissible class, while
repeated compensator-compatible normalization prevents the persistence of an
independent escaping sector. Once both conditions hold, the admissible class is
not merely locally normalized but globally retained under the structural return
mechanisms of the system. This is exactly the content of higher confinement:
the admissible states are no longer only corrected case by case but persist as a
retained confinement regime.
\end{proof}

\begin{lemma}[Non-escape lemma]
\label{lem:non-escape-lemma}
Let $x$ be a closure-stable state in the normalized core. Then no admissible
transition generated within the core sends $x$ into a structurally independent
external sector.
\end{lemma}

\begin{proof}
By closure stability, every admissible normalization step sends $x$ back into
the same closure class. If an admissible transition could produce a genuinely
independent external sector, then closure stability would fail exactly at that
point, because the state would no longer be returnable by the same structural
normalization. Hence admissible transitions may rearrange or re-express states
inside the confinement class, but they cannot export them into an independent
outside sector. This is the structural non-escape property.
\end{proof}

\begin{theorem}[Septinion readability theorem]
\label{thm:septinion-readability-theorem}
If a class is closure-stable and admits the higher confinement reading, then it
is septinion-readable: the resulting stabilized regime may be interpreted as a
septinion/seption-level confinement of the normalized structural body.
\end{theorem}

\begin{proof}
The septinion reading is not introduced as an arbitrary enlargement of the
formalism. Rather, it names the level at which the closure-stable and
non-escaping structural class is most naturally read as a higher unified
container of the admissible sectors. By Lemma~\ref{lem:higher-confinement-lemma},
closure stability already upgrades the system to a higher confinement regime,
and Lemma~\ref{lem:non-escape-lemma} shows that admissible states remain inside
that regime. The septinion/seption reading is therefore the natural structural
reading of this higher confinement container.
\end{proof}

\begin{corollary}[HC as emergent structural consequence]
\label{cor:hc-emergent-structural-consequence}
Within the normalized core, HC is not an externally appended axiom. It appears
as an emergent structural consequence of the one-body, frame-shift, and closure
requirements.
\end{corollary}

\begin{proof}
The generative chain of Phase~I-A forces re-identification and therefore HC,
while the closure chain of Phase~I-B shows that HC is precisely the mechanism
that stabilizes the admissible class. Phase~I-C now identifies the higher
confinement reading selected by this stabilization. Consequently HC is not an
optional addition but the structurally emergent mechanism by which the unified
body remains admissible across shifted readouts and under closure.
\end{proof}

\begin{corollary}[First confinement forcing chain]
\label{cor:first-confinement-forcing-chain}
The first unfolded confinement route of the normalized core is
\[
\text{closure class}
\Longrightarrow
\text{higher confinement}
\Longrightarrow
\text{non-escape}
\Longrightarrow
\text{septinion reading}.
\]
Together with the preceding phases, this yields
\[
\begin{aligned}
&\text{one body} \Longrightarrow \text{re-identification} \Longrightarrow \mathrm{HC}\\
&\Longrightarrow \text{closure} \Longrightarrow \text{higher confinement} \Longrightarrow \text{septinion/seption reading}.
\end{aligned}
\]
\end{corollary}

\begin{proof}
Proposition~\ref{prop:closure-class-proposition} provides the stabilized
closure class. Lemma~\ref{lem:higher-confinement-lemma} upgrades this class to a
higher confinement reading, Lemma~\ref{lem:non-escape-lemma} gives the internal
retention property of admissible transitions, and
Theorem~\ref{thm:septinion-readability-theorem} identifies the corresponding
septinion/seption interpretation. The final displayed chain is obtained by
concatenating this phase with the forcing chains of Phases~I-A and~I-B.
\end{proof}


\subsection*{Phase II-D: Cascade, Readout, and the Collapse-Fixed Sector}
\addcontentsline{toc}{subsection}{Phase II-D: Cascade, Readout, and the Collapse-Fixed Sector}

\begin{lemma}[Cascade admissibility lemma]
\label{lem:cascade-admissibility-lemma}
If $x$ belongs to the closure-stable class selected by the normalized core,
then every admissible cascade image $\hat K(x)$ remains inside the admissible
readout class of the same unified body.
\end{lemma}

\begin{proof}
The cascade is introduced as a composition of projection maps from higher
strata toward the ground sector, not as an ontological splitting of the body.
By Corollary~\ref{cor:first-closure-forcing-chain}, closure-stable states are
precisely those stabilized against drift under shifted readouts. Applying the
cascade to such a state therefore yields another admissible readout of the same
body rather than a new object. Hence admissibility is preserved along the
cascade image of any closure-stable input.
\end{proof}

\begin{lemma}[Terminal ground-state lemma]
\label{lem:terminal-ground-state-lemma}
The ground state $S_0$ is the terminal invariant readout sector of the cascade:
for every admissible state $x$, one has $\hat K(x)=x$ exactly when $x\in S_0$.
\end{lemma}

\begin{proof}
This is the normalized readout form of the ground-state condition already built
into the semantic kernel. The cascade terminates in $S_0$, and terminality
means precisely that no further nontrivial readout step occurs. Therefore the
fixed-point condition $\hat K(x)=x$ is equivalent to membership in the terminal
sector $S_0$.
\end{proof}

\begin{lemma}[Global null/light-axis lemma]
\label{lem:global-light-axis-lemma}
The null-axis structure selected by the normalized core, and read physically as the light axis at stable propagation, is stratified across
the cascade: if a state is read on this null/light axis at one admissible layer, then
its admissible readouts remain aligned with the same global null/light-axis class
through the cascade.
\end{lemma}

\begin{proof}
The Lorentz/null-axis structure, read physically as the light axis, is not introduced as a local afterthought but
as a stratified feature of the whole construction. The cascade only changes the
readout level, not the structural identity of the body. Since admissible
cascade images preserve the same unified body and the same closure-compatible
sector, the null/light-axis reading persists as a global class across the projected
layers.
\end{proof}

\begin{theorem}[Collapse-fixed sector theorem]
\label{thm:collapse-fixed-sector-theorem}
The photon/readout sector identified in the normalized core is the global
collapse-fixed sector of the cascade. In particular, the sector denoted by
$S_0$ is not merely a local null readout but the terminal collapse-fixed sector
selected by the full stratified structure.
\end{theorem}

\begin{proof}
Lemma~\ref{lem:cascade-admissibility-lemma} ensures that the cascade preserves
admissibility, while Lemma~\ref{lem:terminal-ground-state-lemma} identifies
$S_0$ as the terminal invariant sector. Lemma~\ref{lem:global-light-axis-lemma}
shows that the light-axis reading persists globally across admissible layers.
Therefore the photon/readout sector is not obtained from an isolated local
calculation alone; rather, it is the collapse-fixed terminal sector selected by
the entire cascade of admissible projections.
\end{proof}

\begin{lemma}[Readout separation lemma]
\label{lem:readout-separation-lemma}
Within the normalized core one must distinguish three levels: structural
forcing, cascade identification, and physical reading. A statement belongs to
exactly one of these levels unless an explicit bridge is given.
\end{lemma}

\begin{proof}
The structural forcing chain establishes what must hold inside the unified
body. The cascade then identifies how such a structurally forced state appears
at lower readout levels. A physical interpretation is an additional reading of
that identified sector. Conflating these levels obscures where a statement is
proved and where it is interpreted. The separation therefore serves as a
normalization rule for subsequent unfolding steps.
\end{proof}

\begin{corollary}[First readout forcing chain]
\label{cor:first-readout-forcing-chain}
The first unfolded readout route of the normalized core is
\[
\begin{aligned}
&\text{closure-stable class} \Longrightarrow \text{admissible cascade image}\\
&\Longrightarrow \text{terminal ground sector} \Longrightarrow \text{collapse-fixed photon/readout sector}.
\end{aligned}
\]
Taken together with Phases~I-A through~I-C, the normalized core now yields the
combined chain
\[
\begin{aligned}
&\text{one body} \Longrightarrow \text{re-identification} \Longrightarrow \mathrm{HC} \Longrightarrow \text{closure}\\
&\Longrightarrow \text{higher confinement} \Longrightarrow \text{septinion reading}\\
&\Longrightarrow \text{cascade-fixed readout sector}.
\end{aligned}
\]
\end{corollary}

\begin{proof}
The first displayed chain is the direct concatenation of
Lemma~\ref{lem:cascade-admissibility-lemma},
Lemma~\ref{lem:terminal-ground-state-lemma}, and
Theorem~\ref{thm:collapse-fixed-sector-theorem}. The second displayed chain is
obtained by appending this readout route to the forcing chain from
Corollary~\ref{cor:first-confinement-forcing-chain}.
\end{proof}


\subsection*{Phase II-E: Semantic/Tensor Interface and Millennium Intertwining}
\addcontentsline{toc}{subsection}{Phase II-E: Semantic/Tensor Interface and Millennium Intertwining}

\begin{lemma}[Semantic kernel well-formedness lemma]
\label{lem:semantic-kernel-well-formedness-lemma}
The normalized semantic kernel is well formed: its predicates of one-body
identity, triolic organization, orthogonal anchoring, frame-shifted readout,
HC forcing, closure stability, higher confinement, and readout admissibility do
not define competing ontologies but successive structural constraints on one and
the same body.
\end{lemma}

\begin{proof}
By construction, the normalized core begins from the one-body condition and then
adds only constraint predicates that regulate how the same body may be read,
stabilized, and projected. None of the predicates introduces a second carrier
space or an ontologically independent duplicate; each one refines the same
underlying body by imposing an admissibility condition. The semantic kernel is
therefore well formed as a constraint hierarchy rather than a family of
competing state spaces.
\end{proof}

\begin{lemma}[Tensor realization lemma]
\label{lem:tensor-realization-lemma}
The tensorial operators $\mathcal J^\mu{}_{\nu}$, $\mathcal C^\mu{}_{\nu}$,
$\Pi_n{}^\mu{}_{\nu}$, and $\hat K^\mu{}_{\nu}$ realize the normalized semantic
kernel as an operator system: $\mathcal J$ realizes the involutive frame
symmetry, $\mathcal C$ the compensating projection, $\Pi_n$ the cascade steps,
and $\hat K$ the terminal readout map.
\end{lemma}

\begin{proof}
The semantic kernel distinguishes involutive reflection, compensating return,
projection through admissible layers, and terminal ground readout. The tensorial
operators are introduced precisely to encode these roles in compact operator
notation. The semantic and tensorial descriptions therefore stand in a role-wise
correspondence: the operator system is not a second theory but a tensor
realization of the already normalized semantic structure.
\end{proof}

\begin{lemma}[Intertwining lemma]
\label{lem:millennium-intertwining-lemma}
If the Millennium matrix $\mathcal M^\mu{}_{\nu}$ is used as the structural
transport operator between admissible grouped layers, then it must intertwine
with the normalized tensorial kernel. In the minimal normalized form one
requires
\[
\mathcal M\mathcal J=\mathcal J\mathcal M,
\qquad
\mathcal M\mathcal C=\mathcal C\mathcal M,
\qquad
\mathcal M\hat K=\hat K\mathcal M
\]
whenever the corresponding involution, compensation, and cascade-readout
structures are to be preserved across the transport.
\end{lemma}

\begin{proof}
A transport operator that fails to commute with the involution would change the
frame logic of the transported state; one that fails to commute with the
compensator would alter admissibility by moving states in or out of the
compensated sector; and one that fails to commute with the terminal readout map
would not preserve the normalized collapse/readout architecture. Hence a
Millennium transport that is meant to be structure preserving must satisfy the
stated intertwining relations at least for the operators whose semantic role it
claims to preserve.
\end{proof}

\begin{theorem}[Admissibility preservation theorem]
\label{thm:admissibility-preservation-theorem}
If $x$ is admissible in the normalized semantic kernel and $\mathcal M$ is a
Millennium transport satisfying the intertwining relations of
Lemma~\ref{lem:millennium-intertwining-lemma}, then $\mathcal Mx$ remains
admissible in the transported layer.
\end{theorem}

\begin{proof}
Admissibility in the normalized core is carried by compatibility with the one-body
constraint, compensator normalization, and the cascade-compatible readout
structure. By Lemma~\ref{lem:millennium-intertwining-lemma}, the transport
operator preserves the tensorial realizations of exactly these mechanisms.
Therefore applying $\mathcal M$ cannot convert an admissible state into an
inadmissible one merely by transport; it yields an admissible image in the next
layer provided the transport is itself taken within the normalized grouped
architecture.
\end{proof}

\begin{theorem}[Closure preservation theorem]
\label{thm:closure-preservation-theorem}
Under the same hypotheses, closure-stable states are transported by
$\mathcal M$ to closure-stable states of the receiving layer.
\end{theorem}

\begin{proof}
Closure stability is the semantic statement that the admissible state does not
drift out of the normalized class once HC forcing and anchoring are taken into
account. Since admissibility is preserved by
Theorem~\ref{thm:admissibility-preservation-theorem} and the intertwining
relations preserve the compensator and readout structures that stabilize the
closure class, the transported image cannot lose closure merely by layer
transport. Thus closure stability is preserved across admissible Millennium
transport.
\end{proof}

\begin{proposition}[Semantic--tensor equivalence proposition]
\label{prop:semantic-tensor-equivalence-proposition}
The tensorial kernel does not introduce a second structural theory. It is a
parallel notation for the normalized semantic kernel, and the Millennium matrix
is admissible only insofar as its tensorial action represents the same
constraint chain already present on the semantic side.
\end{proposition}

\begin{proof}
By Lemma~\ref{lem:semantic-kernel-well-formedness-lemma}, the semantic side is a
single constrained body. By Lemma~\ref{lem:tensor-realization-lemma}, the
operator side realizes these constraints in tensorial notation. The intertwining
and preservation results show that $\mathcal M$ is not free to create a new
structural doctrine; it is licensed only as a transport operator respecting the
same normalized constraint chain. Hence semantic kernel and tensor kernel are
parallel realizations of one structure.
\end{proof}

\begin{corollary}[First semantic--tensor forcing chain]
\label{cor:first-semantic-tensor-forcing-chain}
The normalized semantic/tensor interface yields the chain
\[
\begin{aligned}
&\text{semantic kernel} \Longrightarrow \text{tensor realization} \Longrightarrow \text{Millennium intertwining}\\
&\Longrightarrow \text{admissibility preservation} \Longrightarrow \text{closure preservation}.
\end{aligned}
\]
\end{corollary}

\begin{proof}
Lemma~\ref{lem:semantic-kernel-well-formedness-lemma} establishes the semantic
kernel, Lemma~\ref{lem:tensor-realization-lemma} gives its tensor realization,
Lemma~\ref{lem:millennium-intertwining-lemma} fixes the admissible transport
relations for $\mathcal M$, and Theorems~\ref{thm:admissibility-preservation-theorem}
and~\ref{thm:closure-preservation-theorem} provide the preservation steps.
\end{proof}


\subsection*{Grouped Layer Architecture and Millennium Transport}
\addcontentsline{toc}{subsection}{Grouped Layer Architecture and Millennium Transport}

The semantic/tensor kernel admits a second normalization in which the document's
many local cores, layers, and readout blocks are collected into three grouped
structural classes. This reduces duplication and makes the role of the
Millennium matrix explicit as a transport operator between admissible groups.

\begin{definition}[Grouped structural architecture]
\label{def:grouped-structural-architecture}
The normalized structural architecture consists of the following three grouped
layers:
\begin{enumerate}[leftmargin=2em,label=(G\arabic*)]
\item \textbf{Generative layer $\mathfrak G$:} the layer of structural formation. It
contains the one-body condition, triolic minimality, orthogonal anchoring,
frame-shifted observation, the ground condition, HC forcing, closure
admissibility, and the higher-confinement/septinion reading.
\item \textbf{Projection/translation layer $\mathfrak P$:} the layer of cascade readout,
sector selection, projection rules, collapse-fixed identifications, null-/light-axis
readings, photon-sector readings, and admissible translation between equivalent
structural descriptions.
\item \textbf{Field/probe realization layer $\mathfrak F$:} the layer of probe states,
two-probe geometry, closure-preserving field equations, interaction rules, and
physical realization of the admissible structure.
\end{enumerate}
These three grouped layers are not independent ontologies. They are different
organizational regimes of the same underlying body and must therefore remain
compatible under admissible transport.
\end{definition}

\begin{definition}[Millennium transport between grouped layers]
\label{def:millennium-transport-grouped-layers}
The Millennium matrix is interpreted at the grouped-layer level as a
structure-preserving converter
\[
\mathcal M_{\mathfrak G\to\mathfrak P}:\mathfrak G\to\mathfrak P,
\qquad
\mathcal M_{\mathfrak P\to\mathfrak F}:\mathfrak P\to\mathfrak F,
\]
and, whenever reverse structural checking is required, also by admissible
back-transport maps
\[
\mathcal M_{\mathfrak P\to\mathfrak G}:\mathfrak P\to\mathfrak G,
\qquad
\mathcal M_{\mathfrak F\to\mathfrak P}:\mathfrak F\to\mathfrak P.
\]
At this level $\mathcal M$ is not merely a notational relabeling but a
converter which must preserve structural admissibility across layer changes.
\end{definition}

\begin{proposition}[Grouped forcing and readout chain]
\label{prop:grouped-forcing-readout-chain}
The grouped architecture organizes the document into the structural chain
\[
\mathfrak G \xrightarrow{\ \mathcal M\ } \mathfrak P
\xrightarrow{\ \mathcal M\ } \mathfrak F,
\]
where:
\begin{enumerate}[leftmargin=2em,label=(\alph*)]
\item $\mathfrak G$ generates the admissible structure,
\item $\mathfrak P$ reads, projects, and translates that structure without changing
its admissibility class,
\item $\mathfrak F$ realizes the same admissible structure as a field/probe system.
\end{enumerate}
In particular, the field/probe layer is not a second origin of the theory; it
is a realization layer of structure first generated in $\mathfrak G$ and made
readable in $\mathfrak P$.
\end{proposition}

\begin{proof}
The generative layer contains the forcing data established above: one-body,
triolic organization, orthogonal anchor, frame shift, HC, and closure. The
projection/translation layer does not create new admissibility; it converts the
generative core into cascade-compatible, sector-resolved, and physically
legible readout. The field/probe layer then realizes those admissible readouts
as explicit probes, field equations, or interaction geometries. Since the same
body is being followed through all three layers, admissible transport by
$\mathcal M$ must preserve rather than replace the structural class.
\end{proof}

\begin{proposition}[Admissibility preservation across grouped layers]
\label{prop:admissibility-preservation-grouped-layers}
Let $x$ be an admissible state generated in $\mathfrak G$. If $\mathcal M$ acts as
an admissible grouped-layer converter, then
\[
\mathrm{Admissible}(x)
\Longrightarrow
\mathrm{Admissible}(\mathcal M_{\mathfrak G\to\mathfrak P}x)
\Longrightarrow
\mathrm{Admissible}(\mathcal M_{\mathfrak P\to\mathfrak F}\mathcal M_{\mathfrak G\to\mathfrak P}x).
\]
Moreover, when closure stability is part of the admissibility data, it is to be
preserved throughout this transport chain.
\end{proposition}

\begin{proof}
By Definition~\ref{def:millennium-transport-grouped-layers}, grouped-layer
transport is only allowed when it preserves structural admissibility. Hence an
admissible state generated in $\mathfrak G$ remains admissible after conversion
into the projection/translation layer, and likewise after conversion into the
field/probe layer. If closure stability were lost during transport, the result
would no longer be a valid readout or realization of the original structure but
a structural rupture. Therefore closure stability must be preserved whenever it
belongs to the originating admissibility class.
\end{proof}

\begin{remark}[Document-normalization consequence]
\label{rem:document-normalization-consequence}
The grouped architecture provides a normalization principle for the text. Local
``cores'' or ``integration blocks'' should be read either as part of
$\mathfrak G$, as part of $\mathfrak P$, or as part of $\mathfrak F$, rather than as
independent foundational centers. This removes duplication: the normative
foundation lies in $\mathfrak G$, the translation/readout logic in $\mathfrak P$,
and the field/probe realization in $\mathfrak F$.
\end{remark}

\subsection*{Phase II-F: Grouped Layer Transport and Ontological Non-Drift}
\addcontentsline{toc}{subsection}{Phase II-F: Grouped Layer Transport and Ontological Non-Drift}

This unfolding block upgrades the grouped architecture from a normalization note
into an explicit transport calculus. The goal is to show that grouped layers do
not introduce new origins of structure, but only different admissible regimes of
one and the same underlying body.

\begin{lemma}[Group decomposition lemma]
\label{lem:group-decomposition}
Every active structural block of the normalized Companion is dominantly
assignable to exactly one of the grouped layers $\mathfrak G$, $\mathfrak P$, or
$\mathfrak F$, even when secondary cross-references to the other layers remain.
\end{lemma}

\begin{proof}
By Definition~\ref{def:grouped-structural-architecture}, the grouped layers are
distinguished by role rather than by separate ontology: generation,
projection/translation, and realization. Any active block therefore has one
primary structural function. Cross-references may connect a block to the other
layers, but they do not erase its dominant role. Hence each active block is
assignable to one grouped layer as its principal location.
\end{proof}

\begin{lemma}[Transport lemma]
\label{lem:group-transport}
Let $x$ be a state generated in $\mathfrak G$. If $\mathcal M_{\mathfrak G\to\mathfrak P}$
is an admissible grouped-layer converter, then its image in $\mathfrak P$ is a
projection/translation of the same admissible structure and not a newly created
structural body.
\end{lemma}

\begin{proof}
By Definition~\ref{def:millennium-transport-grouped-layers}, admissible grouped
transport preserves structural admissibility. Proposition~\ref{prop:grouped-forcing-readout-chain}
states that $\mathfrak P$ reads and translates what $\mathfrak G$ has already
generated. Therefore the image of $x$ under $\mathcal M_{\mathfrak G\to\mathfrak P}$
remains the same admissible structure under a projection/translation regime.
It is not a second origin because admissibility is transported, not recreated.
\end{proof}

\begin{lemma}[Realization lemma]
\label{lem:group-realization}
If $y\in\mathfrak P$ is an admissible projection/translation state and
$\mathcal M_{\mathfrak P\to\mathfrak F}$ is admissible, then its image in
$\mathfrak F$ is a realization of the same admissible structure as a field/probe
state.
\end{lemma}

\begin{proof}
The role of $\mathfrak F$ is realization rather than generation by
Proposition~\ref{prop:grouped-forcing-readout-chain}. Since admissibility is
preserved across grouped transport by Proposition~\ref{prop:admissibility-preservation-grouped-layers},
$\mathcal M_{\mathfrak P\to\mathfrak F}(y)$ remains within the same
admissibility class. Thus the field/probe layer realizes the translated
structure without changing its underlying body.
\end{proof}

\begin{proposition}[No-ontology-drift proposition]
\label{prop:no-ontology-drift}
Admissible transport across the grouped layers does not create ontological
multiplicity. If
\[
x \in \mathfrak G,
\qquad
x_P := \mathcal M_{\mathfrak G\to\mathfrak P}(x),
\qquad
x_F := \mathcal M_{\mathfrak P\to\mathfrak F}(x_P),
\]
then $x$, $x_P$, and $x_F$ are three admissible organizational regimes of one
and the same structural body.
\end{proposition}

\begin{proof}
By Lemma~\ref{lem:group-transport}, passage from $\mathfrak G$ to $\mathfrak P$
produces a translation/readout of the same admissible body. By
Lemma~\ref{lem:group-realization}, passage from $\mathfrak P$ to $\mathfrak F$
produces a field/probe realization of that same admissible body. Since both
steps preserve admissibility rather than replacing it, no new ontology is
introduced along the chain. Therefore $x$, $x_P$, and $x_F$ are structurally
identical in body and differ only by organizational regime.
\end{proof}

\begin{corollary}[First grouped transport forcing chain]
\label{cor:first-grouped-transport-forcing-chain}
The normalized grouped architecture supports the forcing chain
\[
\begin{aligned}
&\mathfrak G \Longrightarrow \text{admissible transport to }\mathfrak P\\
&\Longrightarrow \text{admissible realization in }\mathfrak F \Longrightarrow \text{no ontology drift across grouped layers}.
\end{aligned}
\]
\end{corollary}

\begin{proof}
Combine Lemma~\ref{lem:group-decomposition}, Lemma~\ref{lem:group-transport},
Lemma~\ref{lem:group-realization}, and
Proposition~\ref{prop:no-ontology-drift}.
\end{proof}


\subsection*{Phase III-G: Probe Geometry, Inheritance, and Guided Search}
\addcontentsline{toc}{subsection}{Phase III-G: Probe Geometry, Inheritance, and Guided Search}

This unfolding block upgrades the probe sector from an interpretive add-on to a
structurally inherited realization of the normalized core. The goal is to show
that probes do not introduce a foreign search mechanism, but inherit triolic,
anchored, and closure-compatible structure from the generative and grouped
layers.

\begin{lemma}[Probe inheritance lemma]
\label{lem:probe-inheritance}
If the normalized structural body is triolic, anchored, and closure-admissible,
then every admissible probe realization inherited through $\mathfrak F$ carries a
corresponding triolic and closure-compatible organization.
\end{lemma}

\begin{proof}
By Proposition~\ref{prop:no-ontology-drift}, admissible realization in
$\mathfrak F$ does not create a second ontology, but realizes the same
underlying structural body under a field/probe regime. Hence structural features
that belong to the admissible body---in particular triolic organization and
anchor-compatible closure discipline---remain available in the probe image as
inherited structure rather than externally attached decoration.
\end{proof}

\begin{lemma}[Two-probe compatibility lemma]
\label{lem:two-probe-compatibility}
Let $P_1$ and $P_2$ be two admissible probes directed toward one another. If
both are inherited from the same normalized structural body, then their mutual
configuration is compatible with one common closure-governed geometry.
\end{lemma}

\begin{proof}
By Lemma~\ref{lem:probe-inheritance}, each probe inherits the same triolic and
closure-compatible structural discipline. Since neither probe introduces a new
ontology by Proposition~\ref{prop:no-ontology-drift}, their mutual interaction
must be read inside one common admissibility class. Therefore the two-probe
configuration is not a collision of unrelated structures but a common
closure-governed geometry of one inherited body.
\end{proof}

\begin{lemma}[Orthogonal third-group lemma]
\label{lem:orthogonal-third-group-probe}
In an admissible two-probe configuration, the third structural group remains
orthogonal to both probes and thereby determines the constrained search region.
\end{lemma}

\begin{proof}
The orthogonal third direction is already forced at the generative level by the
orthogonal-anchor principle and remains preserved under admissible realization.
Hence, when two probes are realized as inherited triolic structures, the third
group cannot collapse into either probe direction without violating anchor
compatibility. It therefore remains orthogonal to both and serves as the
stabilizing transversal direction that bounds the admissible search region.
\end{proof}

\begin{proposition}[Probe reduction theorem]
\label{prop:probe-reduction-theorem}
An admissible two-probe configuration reduces to the closure-preserving sector of
the grouped architecture. In particular, the probe interaction may be analyzed
inside the admissible closure class without introducing a second search
ontology.
\end{proposition}

\begin{proof}
By Lemma~\ref{lem:two-probe-compatibility}, the two probes share one common
closure-governed geometry. By Lemma~\ref{lem:orthogonal-third-group-probe}, the
third group fixes the transversal admissible search region. Since closure and
ontology are preserved across grouped realization by
Proposition~\ref{prop:no-ontology-drift}, the probe interaction reduces to the
closure-preserving sector of the same normalized architecture.
\end{proof}

\begin{corollary}[Guided search corollary]
\label{cor:guided-search-corollary}
Auxiliary tools such as frequency extraction, transition guidance, structural
routing, or convergence acceleration act only within the already admissible
probe geometry. They do not generate admissibility, but exploit a search region
that has already been restricted by inheritance, closure, and orthogonal
anchoring.
\end{corollary}

\begin{proof}
By Proposition~\ref{prop:probe-reduction-theorem}, admissible probe analysis is
already confined to the closure-preserving sector. Therefore any auxiliary tool
operates on a domain whose admissibility has been fixed before the tool is
applied. The tool may compress or guide the search, but it does not create the
search geometry itself.
\end{proof}

\begin{corollary}[First probe forcing chain]
\label{cor:first-probe-forcing-chain}
The normalized probe sector supports the forcing chain
\[
\begin{aligned}
&\text{inherited probe structure}
  \Longrightarrow \text{two-probe compatibility}\\
&\Longrightarrow \text{orthogonal third-group search region}\\
&\Longrightarrow \text{closure-preserving probe reduction}\\
&\Longrightarrow \text{guided search inside an already admissible geometry}.
\end{aligned}
\]
\end{corollary}

\begin{proof}
Combine Lemma~\ref{lem:probe-inheritance},
Lemma~\ref{lem:two-probe-compatibility},
Lemma~\ref{lem:orthogonal-third-group-probe},
Proposition~\ref{prop:probe-reduction-theorem}, and
Corollary~\ref{cor:guided-search-corollary}.
\end{proof}



\subsection*{Phase III-H: Structural Master Equation and Exact-Closure Field Regime}
\addcontentsline{toc}{subsection}{Phase III-H: Structural Master Equation and Exact-Closure Field Regime}

This unfolding block upgrades the previously compressed master-equation language
into an explicit field-level forcing chain. The aim is not yet to derive every
physical limit in full detail, but to show that the structural master equation,
its variational reading, the closure-preserving field regime, and the exact-
closure simplifications belong to one continuous argument rather than to loosely
connected interpretive remarks.

\begin{lemma}[Structural master-equation lemma]
\label{lem:structural-master-equation}
The structural master equation is the compressed field-level summary of the
normalized forcing chain from one-body structure through re-identification, HC,
closure, higher confinement, and admissible readout.
\end{lemma}

\begin{proof}
By Corollary~\ref{cor:first-generative-forcing-chain}, the normalized body
forces re-identification and HC. By Corollary~\ref{cor:first-closure-forcing-chain},
this induces a closure class with either exact closure or a controlled residue.
By Corollary~\ref{cor:first-confinement-forcing-chain}, closure yields higher
confinement and the septinion reading. By Corollary~\ref{cor:first-readout-forcing-chain}
and Corollary~\ref{cor:first-semantic-tensor-forcing-chain}, admissible readout
and tensorial transport preserve this structure. Hence the master equation is not
an additional postulate but the compressed field-level expression of the forcing
chain already established.
\end{proof}

\begin{lemma}[Variational consistency lemma]
\label{lem:variational-consistency}
Any variational reading of the structural master equation is admissible only if
it preserves the normalized one-body, HC, and closure conditions. In
particular, admissible variation may unfold the structural equation but may not
replace its forcing chain by an independent ontology.
\end{lemma}

\begin{proof}
By Proposition~\ref{prop:semantic-tensor-equivalence-proposition}, tensorial formulation is
only a parallel realization of the same semantic kernel. Therefore any
variational formulation that is meant to belong to the present architecture must
preserve the same semantic admissibility data. If the variation were to destroy
one-body identity, HC-forcing, or closure preservation, it would cease to be a
variation of the normalized structural equation and would instead define a
second independent system. Hence admissible variation is constrained by the
forcing chain and remains structurally internal to it.
\end{proof}

\begin{lemma}[Closure-field compatibility lemma]
\label{lem:closure-field-compatibility}
The field-level realization of the structural master equation is compatible with
closure precisely when the realized dynamics preserves the HC-selected
admissibility class and does not mix the stable closure sector with an
ontologically foreign complement.
\end{lemma}

\begin{proof}
By Theorem~\ref{thm:closure-preservation-theorem}, admissible transport through
the Millennium interface preserves closure classes. By Proposition~\ref{prop:no-ontology-drift},
realization in the field/probe layer introduces no second ontology. Therefore a
field realization is closure-compatible exactly when it remains inside the
HC-selected admissibility class rather than coupling the closure-stable sector
to a foreign complement. This is the field-level restatement of normalized
closure preservation.
\end{proof}

\begin{theorem}[Exact-closure field theorem]
\label{thm:exact-closure-field-theorem}
In the exactly closed regime, the realized field dynamics reduces to the
closure-stable block of the normalized architecture. Equivalently, the active
field transport becomes block-compatible with the exact closure class.
\end{theorem}

\begin{proof}
By Proposition~\ref{prop:exact-closure-criterion}, exact closure is equivalent
to the disappearance of the nontrivial closure defect. By
Lemma~\ref{lem:closure-field-compatibility}, closure-compatible field dynamics
may then act only on the already stabilized admissibility class. Since no
ontology drift occurs across grouped transport, the realized field dynamics
reduces to the closure-stable block itself. Thus exact closure yields a
block-compatible field regime.
\end{proof}

\begin{lemma}[Minimal compensator disappearance lemma]
\label{lem:minimal-compensator-disappearance}
In the exactly closed regime, the minimal compensator residue vanishes. What
remains is not the disappearance of structural compensation as such, but the
fact that no additional compensator correction is required beyond the exact
closure block.
\end{lemma}

\begin{proof}
By Lemma~\ref{lem:closure-residue-lemma}, the non-exact regime is marked
by a residual deviation from exact closure. By
Theorem~\ref{thm:exact-closure-field-theorem}, exact closure confines the field
transport completely to the closure-stable block. Therefore the additional
minimal compensator term that measures residual off-block correction must
vanish. Structural compensation remains implicit in the exact block, but no
separate compensator residue survives.
\end{proof}

\begin{theorem}[Effective-limit theorem]
\label{thm:effective-limit-theorem}
Any admissible effective limit---for example a geometric, spinorial, or quantum
field-theoretic limit---must arise as a readout of the normalized structural
master equation together with its closure-preserving field realization. Such a
limit may specialize the realization, but cannot bypass the forcing chain that
produces it.
\end{theorem}

\begin{proof}
By Lemma~\ref{lem:structural-master-equation}, the structural master equation is
the compressed summary of the normalized forcing chain. By
Lemma~\ref{lem:variational-consistency} and Lemma~\ref{lem:closure-field-compatibility},
any admissible field realization and any admissible variation remain internal to
that chain. By Theorem~\ref{thm:exact-closure-field-theorem} and
Lemma~\ref{lem:minimal-compensator-disappearance}, the exact regime yields a
pure closure block without additional corrective residue. Hence any effective
limit must be read as a specialization or projection of the normalized
structural field realization rather than as an independent starting point.
\end{proof}

\begin{corollary}[First field forcing chain]
\label{cor:first-field-forcing-chain}
The normalized field sector supports the forcing chain
\[
\begin{aligned}
&\text{structural master equation} \Longrightarrow \text{variational consistency}\\
&\Longrightarrow \text{closure-preserving field compatibility} \Longrightarrow \text{exact closure block regime}\\
&\Longrightarrow \text{vanishing minimal compensator residue} \Longrightarrow \text{admissible effective limit readout}.
\end{aligned}
\]
\end{corollary}

\begin{proof}
Combine the structural master equation, variational consistency, closure-field compatibility, the exact closure-field theorem, minimal compensator disappearance, and the effective limit theorem.
\end{proof}

\subsection*{Concrete mapping of legacy core blocks to grouped layers}
\addcontentsline{toc}{subsection}{Concrete mapping of legacy core blocks to grouped layers}

The grouped normalization can be made explicit by assigning the surviving legacy
core blocks to the three grouped layers $\mathfrak G$, $\mathfrak P$, and
$\mathfrak F$. This turns older version-labelled cores into either active parts
of the normalized architecture or historical/supporting blocks.

\begin{definition}[Layer assignment of legacy core blocks]
\label{def:layer-assignment-legacy-core-blocks}
Within the present monolith, the principal legacy blocks are to be read as
follows.
\begin{enumerate}[leftmargin=2em,label=(L\arabic*)]
\item \textbf{$\mathfrak G$-dominant blocks (generative layer):}
\begin{enumerate}[leftmargin=2.5em,label=(\alph*)]
\item the normalized \emph{Structural Generation Core},
\item \emph{Structural Core and Closure Principles},
\item the generative part of \emph{v1.2.0 Consolidation Core}, especially the
structural reading discipline together with the ground-condition/HC/closure and
higher-confinement logic,
\item the generative part of \emph{v1.4.0 Derivation and Appendix Integration Core},
especially the closure ladder and the micro-to-macro derivation chain.
\end{enumerate}
\item \textbf{$\mathfrak P$-dominant blocks (projection/translation layer):}
\begin{enumerate}[leftmargin=2.5em,label=(\alph*)]
\item \emph{v1.5.0 Projection and Appendix Table Integration Core},
\item the appendix-coupling and projection-summary parts of \emph{v1.4.0 Derivation and Appendix Integration Core},
\item the collapse, light-axis, photon-sector, and cascade-readout parts of the later main sections,
\item all readout, table, projection, and sector-translation interfaces whose
role is to make one admissible structure legible without changing its
admissibility class.
\end{enumerate}
\item \textbf{$\mathfrak F$-dominant blocks (field/probe realization layer):}
\begin{enumerate}[leftmargin=2.5em,label=(\alph*)]
\item the probe-sector and two-probe geometry blocks,
\item the field-equation, interaction, and closure-preserving dynamics sections,
\item the probe-reduction and guided-search realizations insofar as they are
used as operational realizations of previously fixed admissible structure.
\end{enumerate}
\item \textbf{Historical/supporting blocks:} readiness criteria, version-status
passages, and older local consolidation summaries are not to be read as new
foundational centers. They are historical normalization records unless they are
explicitly re-exported into one of $\mathfrak G$, $\mathfrak P$, or
$\mathfrak F$.
\end{enumerate}
\end{definition}

\begin{proposition}[Normalization map for duplicated core material]
\label{prop:normalization-map-duplicated-core-material}
Under the grouped architecture, duplicated structural material is normalized by
the following rule.
\begin{enumerate}[leftmargin=2em,label=(\roman*)]
\item When a legacy block states formation conditions, forcing rules, or
closure-generating principles, it is normalized into $\mathfrak G$.
\item When a legacy block states projection rules, cascade-compatible readout,
sector identifications, appendix transport, or summary-table logic, it is
normalized into $\mathfrak P$.
\item When a legacy block states probes, field equations, interaction channels,
or closure-preserving realizations, it is normalized into $\mathfrak F$.
\item If a block contains mixed material, it is to be read by dominant function,
with the remaining content treated as support, historical context, or material
that should later be split explicitly.
\end{enumerate}
In particular, \emph{v1.2.0 Consolidation Core} and the early derivation cores no
longer define independent foundations beside the normalized structural core;
they contribute either to $\mathfrak G$, to $\mathfrak P$, or to historical
context.
\end{proposition}

\begin{remark}[Immediate practical consequence for the next clean release]
\label{rem:practical-consequence-next-clean-release}
The next clean release should therefore not keep several local cores in
competition. Instead, older version-labelled cores are to be annotated and,
where necessary, shortened according to their dominant grouped-layer function:
formation material belongs to $\mathfrak G$, translation/readout material to
$\mathfrak P$, field/probe realization material to $\mathfrak F$, and pure
status/history material to archival support.
\end{remark}


